\section{Constructing Interfaces: Lattices, Orientations, and Matching Algorithms}\label{sec:constructing_interfaces}

% ---------- %

\subsection{Lattices and Unit Cells}

% ---------- %

In crystalline solids, a \textbf{lattice} is defined as a periodic array of points generated by integer translations of
three fundamental vectors, known as basis vectors (\(\mathbf{a}_1, \mathbf{a}_2, \mathbf{a}_3\)). Mathematically, any
lattice point \(\mathbf{R}\) is expressed as\cite{vasp_bulk}:

% ---------- %

\[
\mathbf{R} = n_1 \mathbf{a}_1 + n_2 \mathbf{a}_2 + n_3 \mathbf{a}_3, \quad n_i \in \mathbb{Z}.
\]

% ---------- %

The choice of lattice vectors (which need not be orthogonal) defines the geometry and symmetry of the
\textit{unit cell}, the smallest repeating structural unit of the crystal.

A given crystal structure can often be described by multiple equivalent representations.

Silicon (Si), for instance, crystallizes in a diamond-cubic structure and can be represented either as:

% ---------- %

\begin{itemize}
    \item A primitive cell containing 2 atoms with non-orthogonal basis vectors, or
    \item A conventional cubic cell containing 8 atoms, exhibiting face-centered cubic symmetry and belonging to space
    \item Group \(Fd\bar{3}m\) (Pearson symbol cF8)\cite{vasp_bulk}.
\end{itemize}

% ---------- %

Both descriptions are equivalent physically, the primitive cell is smaller (fewer atoms, lower symmetry), while the
conventional cell is larger and highlights the cubic symmetry.

% ---------- %

All group IV elements; carbon (C), silicon (Si), germanium (Ge), and alpha-tin (\(\alpha\)-Sn), adopt the diamond-cubic
structure but differ in lattice constants.

For instance, silicon's lattice constant is approximately \(0.5431\,\mathrm{nm}\), whereas germanium's is
around \(0.5658\,\mathrm{nm}\).

These differences create intrinsic lattice mismatches at heterointerfaces, such as a simple \(1{:}1\) Si|Ge interface,
resulting in a lattice mismatch of roughly:

% ---------- %

\[
\epsilon = \frac{a_{\mathrm{Ge}} - a_{\mathrm{Si}}}{a_{\mathrm{Si}}} \approx 4.17\%,
\]

% ---------- %

meaning germanium's lattice is about \(4.2\%\) larger than silicon's.

Bonding these crystals directly forces one side to stretch (under tension) or compress, inducing strain and consequently
increasing the total energy due to lattice misfit.

Significant mismatches can also lead to the formation of misfit defects, such as dislocations, if the strain is not
elastically accommodated\cite{artemis_paper}.
% ---------- %

\subsection{Lattice Mismatch and Supercells}  \label{sec:lattice_mismatching}

% ---------- %

To alleviate strain caused by lattice mismatch at interfaces, one can introduce \textbf{supercells}, enlarged periodic
cells formed by integer expansions \((n_a, n_b, n_c)\) of the original lattice vectors.

A supercell constructed this way contains \(n_a \times n_b \times n_c\) times more atoms than the primitive cell, and
its total energy (in the absence of any interface) scales linearly as:

% ---------- %

\[
E_{\text{supercell}} = E_{\text{unit cell}} \times n_a \times n_b \times n_c.
\]

% ---------- %

For interfaces, the effective mismatch in supercells decreases substantially.

For instance, a Si|Ge \(22{:}21\) supercell reduces lattice mismatch to:

% ---------- %

\[
\epsilon_{\text{eff}} = \frac{22a_{\mathrm{Si}} - 21a_{\mathrm{Ge}}}{22a_{\mathrm{Si}}} \approx 0.016\%,
\]

% ---------- %

resulting in significantly lower strain energy\cite{raffle_paper}.

% ---------- %

In a homogeneous bulk crystal, using a larger supercell does not change the \textit{per-atom} energy (since it's just a
repetition of the same structure).

However, for a heterogeneous interface, choosing an appropriate common supercell can distribute the mismatch strain over
a larger period, thereby reducing the strain per unit length and stabilising the interface.

Essentially, we seek integer multiples of the two lattice constants that coincide as closely as possible.

For example, repeating the Si lattice 25 times and the Ge lattice 24 times yields
$25a_{\mathrm{Si}}\approx 24a_{\mathrm{Ge}}$, leaving an effective mismatch below 0.02\%.

Such a 25:2425{:}24 commensurate supercell drastically minimizes the residual strain energy compared to the original
~4\% mismatch.

In practice, researchers often use this approach (sometimes called a Moiré supercell in 2D contexts) to reduce
mismatches to under ~1\% or less, at which point the interface can be nearly strain-free.

Such a 25:2425{:}24 commensurate supercell drastically minimizes the residual strain energy compared to the original
~4\% mismatch. In practice, researchers often use this approach (sometimes called a Moiré supercell in 2D contexts) to
reduce mismatches to under ~1\% or less, at which point the interface can be nearly strain-free.

% ---------- %

It should be noted that making the supercell very large will reduce strain but at the cost of more atoms in the
simulation.

Thus, there is a trade-off: we want the smallest supercell that brings the lattices into acceptable registry.

As the supercell size grows, the interfacial energy (per area) will converge to a stable value because the misfit strain
is distributed over more atoms, and the system can relax closer to bulk-like conditions.

% ---------- %

\subsection{Orientation and Planar Matching}

% ---------- %

Besides matching lattice constants, crystallographic orientation plays a crucial role in constructing realistic
interfaces.

Crystal orientations are commonly described by Miller indices \((hkl)\), which specify the orientation of a
crystallographic plane relative to the unit cell axes.

For example, a (100) surface represents a cut perpendicular to the \(x\)-axis of a cubic cell, a (110) surface is a
diagonal cut, and a (111) plane slices diagonally across the cubic cell.

When constructing an interface, one must carefully select which Miller-indexed plane of material A will contact a
particular plane of material B. Rotating and aligning slabs according to these indices determine the atomic registry at
the interface, significantly affecting interfacial properties such as bonding patterns, strain distribution, and
energetics\cite{artemis_paper}.

% ---------- %

Different Miller planes offer distinct surface terminations, each with unique atomic configurations.

Thus, even a single crystallographic orientation may yield several energetically and structurally distinct interfaces
depending on the chosen termination.

Moreover, some seemingly mismatched orientations can be rendered commensurate by careful in-plane rotation.

For instance, high-throughput searches have identified viable alignments between dissimilar crystal planes (e.g.
cubic (010) against hexagonal-like (111)), minimising residual mismatch through rotation.

% ---------- %

Additionally, crystals with differing stacking sequences can form interfaces: a face-centered cubic (FCC) (111) surface
has ABC atomic stacking, whereas a hexagonal close-packed (HCP) (0001) surface follows ABAB stacking.

Joining these surfaces can lead to multiple possible stacking registries (atomic alignments across the interface), each
with unique structural and energetic characteristics.

% ---------- %

Interface registry involves two primary degrees of freedom:

% ---------- %

\begin{itemize}
    \item \textit{In-plane shifts} (\(\mathbf{d}_{\parallel}\)): lateral translations within the interface plane,
    altering atomic overlap and influencing bonding arrangements.
    \item \textit{Out-of-plane separations} (\(\mathbf{d}_{\perp}\)): vertical adjustments controlling bonding distances
    between slabs, or introducing vacuum gaps to simulate weakly bonded or non-bonded interfaces.
\end{itemize}

% ---------- %

Together, these parameters fully define the interfacial configuration space, requiring systematic exploration to
identify optimal interface structures.

% ---------- %

\subsection{Automated Interface Construction (ARTEMIS)}

% ---------- %

To systematically explore the large configurational space involved in constructing interfaces and identify physically
plausible, low-strain structures, automated algorithms like ARTEMIS are employed\cite{artemis_paper}.

ARTEMIS automates the construction of chemically realistic interfaces by following a structured, step-by-step workflow
(\figref{fig:artemis_uml}):

% ---------- %

\begin{enumerate}
    \item \textbf{Surface selection:} The user specifies Miller indices \((hkl)_A, (h'k'l')_B\) for surfaces of
    materials A and B, typically choosing low-energy or experimentally relevant facets.
    \item \textbf{Slab generation:} Each bulk crystal is cleaved along the chosen Miller planes, producing slabs with
    sufficient vacuum to avoid interactions between periodic images and ensuring stoichiometric and chemically
    reasonable terminations. Multiple possible surface terminations can be considered if relevant.
    \item \textbf{2D lattice matching:} ARTEMIS employs a lattice-matching algorithm based on the Zur-McGill method,
    seeking integer matrices \(M_A\) and \(M_B\) that closely approximate the two materials' lattice vectors:
    \[
    M_A [\mathbf{a}_1^A, \mathbf{a}_2^A] \approx M_B [\mathbf{a}_1^B, \mathbf{a}_2^B],
    \]
    thereby minimizing lattice mismatch strain\cite{artemis_paper}. Candidate supercells (integer multiples of slab
    lattice vectors) are generated and compared to find commensurate matches within a small strain tolerance (typically
    \(<5\%\)). Slight rotations can also be explored at this stage to further reduce mismatch.
    \item \textbf{Registry enumeration (in-plane shifts):} Given a commensurate lattice match, ARTEMIS systematically
    explores multiple lateral translations (\(\mathbf{d}_{\parallel}\)) of one slab relative to the other.

    Fractional offsets (grid sampling) ensure that all high-symmetry and distinct stacking configurations are sampled,
    significantly influencing the resulting atomic registry and bonding environment.
    \item \textbf{Vertical optimization (out-of-plane separation):} For each in-plane shifted configuration, the
    perpendicular slab separation (\(\mathbf{d}_{\perp}\)) is adjusted to ensure realistic bonding distances.

    Slabs are moved toward each other until physically meaningful interatomic distances or bonding criteria are
    satisfied, avoiding atomic overlaps or excessive gaps.
    \item \textbf{Filtering candidates:} To manage computational complexity, ARTEMIS filters interface candidates based
    on:
        \begin{itemize}
            \item Maximum allowable supercell size (limiting the total atom count).
            \item Maximum strain thresholds (usually below 5\%) to ensure structural plausibility.
            \item Chemical feasibility, such as minimizing dangling or unsatisfied bonds at the interface.
        \end{itemize}
    \item \textbf{Output generation:} Finally, ARTEMIS outputs the most feasible interface candidates as initial
    structures for further relaxation and analysis via density functional theory (DFT) or machine-learned potentials
    (MLPs).

    Outputs include metadata such as applied strain, shift vectors, and Miller indices to facilitate subsequent
    identification and tracking.
\end{enumerate}

% ---------- %

\begin{figure}[H]
    \centering
    \includegraphics[width=0.85\textwidth]{figures/artemis_algorithm.png}
    \caption{UML-style flowchart detailing ARTEMIS algorithm stages: slab generation, Miller-plane matching, supercell
    optimization, registry sampling, and filtering. Adapted from Taylor et al.\cite{artemis_paper}.}
    \label{fig:artemis_uml}
\end{figure}

% ---------- %

By automating these steps, ARTEMIS streamlines the otherwise laborious and complex task of interface generation,
enabling systematic, unbiased exploration of potential heterostructure interfaces.

This method has successfully modeled diverse interfaces ranging from semiconductor heterojunctions to 2D/3D material
contacts, significantly accelerating computational materials research and enhancing the reliability of interface
predictions\cite{artemis_paper,raffle_paper}.

% ---------- %

\label{para:polymorph_interface_reconstruction} % ---------- %

When the interface forms between distinct structural phases, such as cubic and hexagonal polymorphs, unique bonding
geometries and strain fields arise\cite{Chagarov2015-ol, Di_Liberto2022-gz}.

For instance, in layered 2D|3D systems like \ce{HfS2}|\ce{HfO2}, the band alignment changes depending on whether the
materials are stacked or laterally connected, with lateral interfaces showing increased reconstruction due to bonding
mismatch between the different crystalline motifs\cite{Taylor2023-wh}.

These reconstructions are not merely geometric; they result in different electrostatic and electronic behaviours,
influencing band offsets and charge transfer dynamics.

\label{para:interfacial_defects} % ---------- %

Defects, including vacancies, interstitials, dislocations, and grain boundaries, are often intrinsic to interface
formation and can profoundly modify interfacial properties\cite{Maji2021-kc}.

They may serve as scattering centres, recombination sites, or sources of local doping\cite{Leitherer2023-ew}.

In grain boundaries of polycrystalline materials, for example, variations in local stoichiometry and strain can create
regions of localised metallicity or enhanced dielectric response, as shown in systems exhibiting colossal permittivity\cite{}.

In some systems, such as CCTO (\ce{CaCu3Ti4O12}), interfacial defects and microstructural features have been shown to
dominate the macroscopic dielectric properties\cite{RAFFLE}.

\label{para:defect_ml_prediction} % ---------- %

In computational studies, predicting the impact of such defects requires careful treatment of their spatial distribution,
charge state, and interaction with surrounding lattice atoms\cite{}.

Modern approaches like the RAFFLE algorithm allow for machine learning-driven structure prediction that accounts for
void-filling and local bonding environments, enabling more realistic models of defect-laden interfaces\cite{RAFFLE}.

\label{para:interface_applications} % ---------- %

Depending on their nature, interfaces may enhance or impede material performance\cite{Leitherer2023-ew, Butler2019-ws}.

This subsection surveys several application domains in which understanding and engineering interfacial properties is
essential.

\label{para:gate_dielectric_interface} % ---------- %

In semiconductor transistor devices, the interface between the gate dielectric and the semiconducting channel is a
critical design parameter\cite{Schusteritsch2015-lw, Wild2025-jy}.

% ---------- %

As traditional \ce{SiO2} dielectrics are replaced by high-$\kappa$ materials such as \ce{HfO2} to meet miniaturisation
demands, understanding charge transfer and band alignment across interfaces like \ce{HfS2}| \ce{HfO2} becomes essential
for ensuring device reliability and performance\cite{}.

Notably, in such 2D|3D heterostructures, the type and strength of bonding (e.g. van der Waals vs. chemical) directly
influences charge redistribution and can cause significant deviation from band alignments predicted by Anderson's rule,
the vacuum levels of the two materials must be aligned\cite{Davies2021-zh}.

\label{para:tmdc_interface_band_alignment} % ---------- %

In layered 2D electronics, especially those using TMDCs, the stacking geometry, layer
orientation, and interfacial disorder all strongly influence the band alignment\cite{Davies2024-ju}.

These factors affect carrier dynamics, optical transitions, and switching behaviour.

Traditional models such as Anderson's rule have proven insufficient; corrections using terms such as $\Delta E_\Gamma$ and
$\Delta E_{\text{IF}}$ have been developed to capture hybridisation and interface dipole contributions more accurately\cite{Butler2019-ws}.

% todo: add formula for detla E_\Gamma and delta E_IF

These corrections allow for predictive modelling of heterostructure behaviour and support targeted band offset
engineering.

\label{para:energy_storage_interface} % ---------- %

Interfaces also govern performance in energy storage systems\cite{Leitherer2023-ew}.

In all-solid-state batteries, for example, the interface between electrode and solid electrolyte dictates ion mobility,
chemical stability, and space charge layer formation\cite{Leung2020-tv, Butler2019-ws}.

Such interfacial processes can lead to resistive interphases or dendrite formation, reducing battery life and safety\cite{Butler2019-ws, Leung2020-tv}.

Density Functional Theory (DFT) modelling has shown that interfacial reconstructions and electrostatic gradients can
dominate device behaviour\cite{Butler2019-ws, Leung2020-tv}.

\label{para:thermoelectric_interface} % ---------- %

In thermoelectric materials, interfaces are employed to decouple electrical and thermal transport\cite{Leitherer2023-ew}.

Interface-induced phonon scattering suppresses thermal conductivity while preserving electronic conductivity; enhancing
the thermoelectric figure of merit, $ZT$\cite{Hepplestone2010-vw}.

Layered TMDCs and heterostructures with nanostructured interfaces exploit this mechanism to improve thermal management\cite{Low2025-on}.

\label{para:interfaces_active_modelling_challenge} % ---------- %

Across these examples, it is evident that interfaces are not passive boundaries, but active, tunable regions that
critically determine macroscopic functionality\cite{Leitherer2023-ew, Di_Liberto2022-gz}.

Their predictive modelling, through DFT, machine learning potentials, or interface-specific structure search tools like
ARTEMIS, remains a central challenge in modern materials design\cite{}.

% ---------- %
